% Bài giảng Toán 8 - Hằng đẳng thức
\section{Hằng đẳng thức đáng nhớ: $(a+b)^2 = a^2 + 2ab + b^2$ và $\sqrt{a^2+b^2}$}

\subsection{I. Lý thuyết trọng tâm}
Trong toán học đại số lớp 8, các hằng đẳng thức đáng nhớ là công cụ cơ bản để biến đổi biểu thức đại số.

\begin{align*}
(a + b)^2 &= a^2 + 2ab + b^2 \\
(a - b)^2 &= a^2 - 2ab + b^2 \\
a^2 - b^2 &= (a - b)(a + b)
\end{align*}

Bảng đối chiếu một số biểu thức cụ thể:
\begin{center}
\begin{tabular}{|c|c|c|}
\hline
Biểu thức $A$ & Biểu thức $B$ & Khai triển $(A+B)^2$ \\
\hline
$x$ & $1$ & $x^2 + 2x + 1$ \\
$2x$ & $3$ & $4x^2 + 12x + 9$ \\
\hline
\end{tabular}
\end{center}

\subsection{II. Các dạng toán cơ bản}

\subsubsection{Dạng 1: Khai triển biểu thức}
\begin{vd}
Khai triển biểu thức $P = (2x + 3y)^2$.
\loigiai{
Ta áp dụng hằng đẳng thức bình phương một tổng:
$P = (2x)^2 + 2 \cdot (2x) \cdot (3y) + (3y)^2 = 4x^2 + 12xy + 9y^2$.
}
\end{vd}

\subsubsection{Dạng 2: Rút gọn biểu thức}
\begin{bt}
Rút gọn biểu thức $Q = (x + 1)^2 - (x - 1)^2$.
\loigiai{
Khai triển từng hằng đẳng thức:
$Q = (x^2 + 2x + 1) - (x^2 - 2x + 1) = 4x$.
}
\end{bt}

\subsection{III. Bài tập trắc nghiệm tự luyện}

\begin{ex}
Khai triển hằng đẳng thức $(x - 2)^2$ ta được kết quả nào sau đây?
\choice
{$x^2 - 4x + 4$}
{$x^2 - 2x + 4$}
{$x^2 + 4x + 4$}
{$x^2 - 4$}
\loigiai{
Áp dụng công thức $(a - b)^2 = a^2 - 2ab + b^2$ với $a = x$ và $b = 2$, ta được $x^2 - 4x + 4$.
}
\end{ex}
